%% kapa: exam=1 round=27 year=2016 period=1 form=A subject=real_estate subject_ko=부동산학원론 seq=27 number=107
%% source: Q-Net 감정평가사 기출문제 (cst00309) artlSeq=5169402 file=extracted/1008/5169402/제27회 감정평가사 1차 문제/제27회 감정평가사 1교시문제(A형).pdf
%% conversion: kordoc 4.24.51 -> tools/gen1.py (auto); review=proofread-sonnet; flags=none
\begin{question}{107}
다음은 부동산투자의 예상 현금흐름표이다. 이 투자안의 수익성지수(PI)는? (단, 현금유출은 기초, 현금유입은 기말로 가정하고, 0년차 현금흐름은 현금유출이며, 1년차부터 3년차까지의 현금흐름은 연 단위의 현금유입만 발생함. 할인율은 연 10 \%이고, 주어진 조건에 한함)
\begin{flushright}(단위: 만원)\end{flushright}
{\small\centering
\begin{tabular}{|c|c|c|c|c|}\hline
사업 기간 & 0년 & 1년 & 2년 & 3년 \\ \hline
현금흐름 & 2,000 & 550 & 1,210 & 1,331 \\ \hline
\end{tabular}\par}
\begin{choices}
  \item 1.15
  \item 1.25
  \item 1.35
  \item 1.40
  \item 1.45
\end{choices}
\end{question}
\answer{2}
